% !TEX TS-program = pdflatex
% !TEX encoding = UTF-8 Unicode

% This file is a template using the "beamer" package to create slides for a talk or presentation
% - Giving a talk on some subject.
% - The talk is between 15min and 45min long.
% - Style is ornate.

% MODIFIED by Jonathan Kew, 2008-07-06
% The header comments and encoding in this file were modified for inclusion with TeXworks.
% The content is otherwise unchanged from the original distributed with the beamer package.

\documentclass{beamer}
\setbeamercovered{invisible}



% Copyright 2004 by Till Tantau <tantau@users.sourceforge.net>.
%
% In principle, this file can be redistributed and/or modified under
% the terms of the GNU Public License, version 2.
%
% However, this file is supposed to be a template to be modified
% for your own needs. For this reason, if you use this file as a
% template and not specifically distribute it as part of a another
% package/program, I grant the extra permission to freely copy and
% modify this file as you see fit and even to delete this copyright
% notice. 


\mode<presentation>
{
  \usetheme{Warsaw}
  % or ...

  % or whatever (possibly just delete it)
}


\usepackage[english]{babel}
% or whatever

\usepackage[utf8]{inputenc}
% or whatever

\usepackage{times}
\usepackage[T1]{fontenc}
% Or whatever. Note that the encoding and the font should match. If T1
% does not look nice, try deleting the line with the fontenc.

%%% MATH RELATED

\input{macros.tex}

\title{MATH 350-2 Advanced Calculus} 
\subtitle
{} % (optional)

\author[W.R. Casper] % (optional, use only with lots of authors)
{W.R. Casper}
% - Use the \inst{?} command only if the authors have different
%   affiliation.

\institute[California State University Fullerton] % (optional, but mostly needed)
{
  Department of Mathematics\\
  California State University Fullerton}
% - Use the \inst command only if there are several affiliations.
% - Keep it simple, no one is interested in your street address.

\subject{Talks}
% This is only inserted into the PDF information catalog. Can be left
% out. 



% If you have a file called "university-logo-filename.xxx", where xxx
% is a graphic format that can be processed by latex or pdflatex,
% resp., then you can add a logo as follows:

% \pgfdeclareimage[height=0.5cm]{university-logo}{university-logo-filename}
% \logo{\pgfuseimage{university-logo}}



% Delete this, if you do not want the table of contents to pop up at
% the beginning of each subsection:
\AtBeginSubsection[]
{
  \begin{frame}<beamer>{Outline}
    \tableofcontents[currentsection,currentsubsection]
  \end{frame}
}


% If you wish to uncover everything in a step-wise fashion, uncomment
% the following command: 

%\beamerdefaultoverlayspecification{<+->}


\begin{document}

\begin{frame}
  \titlepage
\end{frame}

\begin{frame}{Outline}
  \tableofcontents
  % You might wish to add the option [pausesections]
\end{frame}

% Since this a solution template for a generic talk, very little can
% be said about how it should be structured. However, the talk length
% of between 15min and 45min and the theme suggest that you stick to
% the following rules:  

% - Exactly two or three sections (other than the summary).
% - At *most* three subsections per section.
% - Talk about 30s to 2min per frame. So there should be between about
%   15 and 30 frames, all told.

\section{Real Analysis Lecture 21}

\subsection{Infinite series}

\begin{frame}{Basic definition}
\begin{defn}
An \textbf{infinite series} is an ordererd pair of sequences $(\{a_n\},\{s_n\})$ where 
$$s_n = a_1+a_2+\dots+a_n.$$
The sequence $\{s_n\}$ is called the \textbf{sequence of partial sums} and $s_n$ is called the $n$'th partial sum.
We say the series converges or diverges if $\{s_n\}$ converges or diverges.
\end{defn}
\textbf{Notation:}
$$a_1+a_2+a_3+\dots\quad\text{or}\quad \sum_{n=1}^\infty a_n.$$
\end{frame}

\begin{frame}{Basic definition}
We write
$$\sum_{n=1}^\infty a_n = L$$
to mean that the limit of the sequence $s_n$ is $L$.
We call $L$ the \textbf{sum} of the series.
\end{frame}

\begin{frame}{Examples}
$$\sum_{n=1}^\infty \frac{1}{n^2+n} = 1.$$
\pause
Telescoping:
\pause
\begin{align*}
s_n &= \frac{1}{1^2+1} + \frac{1}{2^2+2} +  \frac{1}{3^2+3} + \dots + \frac{1}{n^2+n}\\
    &= \left(\frac{1}{1}-\frac{1}{2}\right) + \left(\frac{1}{2}-\frac{1}{3}\right) +\dots+ \left(\frac{1}{n}-\frac{1}{n+1}\right)\\
    &= 1-\frac{1}{n+1}.
\end{align*}
\end{frame}

\begin{frame}{Examples}
$$\sum_{n=1}^\infty r^n = \frac{r}{1-r},\quad |r| < 1.$$
\pause
Geometric:
\pause
$$s_n = r + r^2 + r^3 + \dots + r^n= r(1+r+r^2+\dots+r^{n-1}) = r(1+s_{n-1})$$
\pause
$$s_n = r + r^2 + r^3 + \dots + r^n= (r+r^2+\dots+r^{n-1}) + r^n = s_{n-1} + r^n$$
\pause
$$s_n = \frac{r-r^{n+1}}{1-r}$$
\end{frame}

\begin{frame}{Examples}
$$\sum_{n=1}^\infty (-1)^{n+1}\frac{1}{n} = \ln(2)$$
\pause
Taylor series:
\pause
$$\ln(1-x) = -x - \frac{1}{2}x^2 - \frac{1}{3}x^3 - \frac{1}{4}x^4 - \dots$$
\pause
Converges for $-1\leq x < 1$.
\end{frame}

\begin{frame}{Tests for Convergence}
\begin{thm}[Divergence Test]
If $\lim_{n\rightarrow\infty} a_n$ either does not exist or is equal to something nonzero, then the series $\sum a_n$ diverges.
\end{thm}
\end{frame}

\begin{frame}{Tests for Convergence}
\begin{thm}[Cauchy condition]
The series $\sum a_n$ converges if and only if for all $\epsilon > 0$ there exists $N\in\mathbb Z_+$ such that $n\geq N$ implies
$$\left\lvert a_{n} + a_{n+1} + \dots + a_{n+p}\right\rvert < \epsilon$$
for all $p\in\mathbb Z_+$.
\end{thm}
\end{frame}

\begin{frame}{Tests for Convergence}
\begin{defn}
An \textbf{alternating series} is a series of the form $\sum (-1)^{n+1}a_n$ with $a_n\geq 0$ for all $n$.
\end{defn}
\begin{thm}[Alternating Series Test]
If $a_n$ is a decreasing sequence converging to $0$, then $\sum (-1)^{n+1}a_n$ converges to a real number $s$ and
$$|s-s_n| < a_{n+1}\quad\forall\ n\in\mathbb Z_+.$$
\end{thm}
\end{frame}

\begin{frame}{Tests for Convergence}
\begin{defn}
We say $\sum a_n$ is \textbf{absolutely convergent} if $\sum |a_n|$ converges.\\
We say $\sum a_n$ is \textbf{conditionally convergent} if it is convergent but not absolutely convergent.
\end{defn}
\begin{thm}[Absolute Convergence Test]
If a series is absolutely convergent, then it is convergent.
\end{thm}
\end{frame}

\end{document}


